% Generated draft: requires mathematical and source review.
\documentclass[11pt]{article}
\usepackage[margin=25mm]{geometry}
\usepackage{fontspec}
\setmainfont{Latin Modern Roman}
\usepackage{amsmath,amssymb,mathtools,unicode-math}
\setmathfont{Latin Modern Math}
\usepackage{enumitem,needspace,xurl,hyperref}
\setlength{\parindent}{0pt}
\setlength{\parskip}{5pt}
\setlength{\emergencystretch}{3em}
\newcommand{\N}{\mathbb{N}}
\newcommand{\Z}{\mathbb{Z}}
\newcommand{\Q}{\mathbb{Q}}
\newcommand{\R}{\mathbb{R}}
\newcommand{\C}{\mathbb{C}}
\newcommand{\F}{\mathbb{F}}
\newcommand{\A}{\mathbb{A}}
\newcommand{\PP}{\mathbb{P}}
\newcommand{\E}{\mathbb{E}}
\newcommand{\eps}{\varepsilon}
\newcommand{\dd}{\,\mathrm{d}}
\begin{document}
\section*{UnsolvedMath ID: 700641}
\textbf{Source problem code:} catalogue entry
\subsection{Definitions and mathematical statement}
Let \(F\) be a free group with identity \(1_F\). For \(\varphi\in\operatorname{Aut}(F)\) and \(g\in F\), write \(g^\varphi=\varphi(g)\). Let \(x_1,\ldots,x_m\) denote the distinct unknowns occurring in the equation, and let \(i_1,\ldots,i_n\in\{1,\ldots,m\}\); indices may repeat, so different automorphisms may be applied to the same unknown.

\textbf{Part (b): requested decision problem.} Given a word expression \(w\) in \(n\) argument slots and automorphisms \(\varphi_1,\ldots,\varphi_n\in\operatorname{Aut}(F)\), does there exist an algorithm that always halts and decides whether
\[
\exists(g_1,\ldots,g_m)\in F^m\qquad
w\bigl(\varphi_1(g_{i_1}),\ldots,\varphi_n(g_{i_n})\bigr)=1_F?
\]
Here the automorphisms are prescribed data, and the elements \(g_1,\ldots,g_m\) are the unknowns. For a coefficient-free group word, \(w\) is a finite word in formal symbols \(Y_1^{\pm1},\ldots,Y_n^{\pm1}\), evaluated by substitution. The supplied record does not specify whether this coefficient-free convention is intended or whether coefficients from \(F\) are permitted. It also does not specify the rank of \(F\), whether \(F\) is fixed or varies with the input, or effective representations of the group and automorphisms. Thus the displayed solvability predicate is explicit, but the precise scope of the algorithmic question remains qualified by these unresolved input conventions.

\textbf{Archived part (a).} Given \(a,b\in F\) and \(\varphi\in\operatorname{Aut}(F)\), does there exist an algorithm deciding whether
\[
\exists x\in F\qquad a\varphi(x)=xb?
\]
The supplied archive reports an affirmative answer to part (a) and cites the references listed below. The original text suggests that this question may be useful for solving the problem of equivalence of two knots; this is preserved as source motivation, without asserting a reduction or application. The supplied registry labels the overall problem partially solved, with part (a) archived and part (b) retained in the main section. These are supplied source classifications, not an independently verified assessment.
\subsection{Short English statement}
Part (b) asks for an algorithm deciding whether a word equation has a solution in a free group when each occurrence of an unknown may be transformed by a prescribed automorphism. Repeated indices refer to the same unknown. The intended free-group rank, input representations, uniformity, and allowance of coefficients need clarification. Part (a), concerning \(a\varphi(x)=xb\), is reported solved in the supplied archive.
\subsection{Sources}
\begin{itemize}
\item Supplied local UnsolvedMath registry record: permanent ID \(700641\), exact source problem code \texttt{catalogue entry}, title \emph{Kourovka Notebook Problem 10.26}; upstream revision \texttt{d371ee917ac9647fbb06cd6fc529deec6f3dda82}. Proposed by G. S. Makanin in 1986.
\item \emph{The Kourovka Notebook (Unsolved Problems in Group Theory)}, 10th issue (1986), Problem 10.26; supplied record describes an update in the 21st issue (2026), \texttt{arXiv:1401.0300v48}. The registry describes the collection as Issues 1--21 (1965--2026), edited by E. I. Khukhro and V. D. Mazurov. This attribution is supplied and has not been independently checked.
\item Archive of Solved Problems reference for the affirmative answer to part (a): O. Bogopolski, A. Martino, O. Maslakova, E. Ventura, \emph{Bull. London Math. Soc.} \textbf{38}, no. 5 (2006), 787--794. The supplied record gives no article title or URL.
\item Further generalizations and applications are attributed by the supplied archive to O. Bogopolski, A. Martino, E. Ventura, \emph{Trans. Amer. Math. Soc.} \textbf{362} (2010), 2003--2036. The supplied record gives no article title or URL and does not state that this reference resolves part (b).
\item Supplied status review dated \texttt{2026-10-16}: canonical status \texttt{partially\_solved}, evidence labelled \texttt{published}; its notes place solved part (a) in the archive and remaining part(s) in the main section. This is supplied dated curation, not an independently verified literature assessment. The status-review note truncates its second bibliographic reference; the fuller reference above is preserved from the supplied background.
\end{itemize}
\end{document}
